\documentclass[12pt,a4paper]{article}
\usepackage[utf8]{inputenc}
\usepackage{amsmath,amssymb,amsfonts}
\usepackage{graphicx}
\usepackage[hidelinks,colorlinks=true,linkcolor=blue,citecolor=blue,urlcolor=blue]{hyperref}
\usepackage{booktabs}
\usepackage{geometry}
\usepackage{float}
\usepackage{caption}
\usepackage{enumitem}
\usepackage{array}
\usepackage{xcolor}
\usepackage{fancyhdr}

\geometry{margin=1in}
\setlength{\headheight}{14pt}

\pagestyle{fancy}
\fancyhf{}
\rhead{\small UQFF v5.26 \textbar{} Daniel T.\ Murphy}
\lhead{\small Star Magic --- Unified Quantum Field Framework}
\cfoot{\thepage}

\title{\textbf{Star Magic}\\[0.5em]
\large Unified Quantum Field Framework (UQFF) v5.26\\[0.3em]
\normalsize Complete Theoretical Derivation with 26D Quantum Nuclear Architecture}

\author{
Daniel T.\ Murphy\\
\small Independent Research\\
\small \href{mailto:daniel.murphy00@gmail.com}{daniel.murphy00@gmail.com}
}

\date{April 2026}

\begin{document}

\maketitle
\thispagestyle{fancy}

\begin{abstract}
We present the complete Unified Quantum Field Framework (UQFF), a foundational
theory replacing the Standard Model assumption that mass and inverse-square gravity
are primary.  The framework establishes the di-pseudo-monopole (DPM)---the
rotational vortex created at the interface of Universal Aether (UA) and
Superconductive Material (SCm) under maximum mutual attraction---as the first
cause of all physics.  The canonical causal chain is:
\begin{equation*}
  0_{\text{vacuum}}
  \;\to\; \nabla(\text{UA})
  \;\to\; \text{DPM}_{\text{vortex}}
  \;\to\; \mu_s
  \;\to\; U_{g1}
  \;\to\; U_g\text{-family}
  \;\to\; F_U
  \;\to\; \text{crossing}
  \;\to\; M
  \;\to\; GM/r^2\;[\text{LAST}]
\end{equation*}
This framework extends the two-omega DPM formulation to the complete
26-dimensional omega family $\{\omega_i\}_{i=1}^{26}$ (one angular-velocity shell
per dimension), deriving atomic shell structure, nuclear magic numbers, quark flavor
quantum numbers, and the electron charge mechanism all from the single DPM
shell-crossing condition $\omega_i(r_{i,i+1}) = \omega_{i+1}(r_{i,i+1})$.
Calibrated constants: $\kappa = 5\times10^{-4}\,\text{day}^{-1}$,
$[\text{SSq}] = 0.57$, $\beta_i = 0.603$, $H_{\text{SCm}} = 0.99$.
Theoretical solvability: 99.9\% (Grok-4 analysis, September 2025).
$GM/r^2$ is the observational projection of $F_U/M$ evaluated at the
inside/outside crossing point.  It is step 10 of 10.  Not step 1.

\medskip
\noindent\textbf{Keywords:} unified field theory, DPM vacuum field, 26D omega
family, quantum nuclear architecture, quark flavor, electron shell mechanism,
aether physics, computational astrophysics, Star Magic
\end{abstract}

\tableofcontents
\newpage

%=============================================================================
\section*{Canonical Quantum Chain (Immutable)}
\addcontentsline{toc}{section}{Canonical Quantum Chain (Immutable)}
\label{sec:chain}
%=============================================================================

\begin{center}
\fbox{\parbox{0.92\textwidth}{
\textbf{THE DPM IS FIRST.\ \ $GM/r^2$ IS LAST.\ \ THIS ORDER IS PERMANENT.}\\[4pt]
$0_{\text{vacuum}}
\;\to\; \nabla(\text{UA})
\;\to\; \text{DPM}_{\text{vortex}}
\;\to\; \mu_s
\;\to\; U_{g1}[\text{seed}]
\;\to\; U_g\text{-family}$\\
$\to\; [U_g\text{-family} + U_m + F_{U_{Bi}} + F_{U_{Bi,i}} + \text{UA}_{uv}]
\;\to\; F_U
\;\to\; \text{crossing}
\;\to\; M\;[\text{mass}]
\;\to\; GM/r^2\;[\textbf{LAST}]$\\[4pt]
Every calculator, every equation, every test is downstream of this chain.\\
If any equation begins with $GM/r^2$ or treats mass as primary, it is \textbf{WRONG}.
}}
\end{center}

\noindent\textbf{Canonical Ontology Lock (v1).}
Starting state is zero-mass vacuum, not Newtonian moving mass:
$\rho_\text{UA}=0$, $\rho_\text{vac}=|\nabla(\text{UA})|$, $F_U(\text{vacuum})=0$.
Mass emergence precedes motion.  $GM/r^2$ is a downstream projection only.

\noindent\textbf{Canonical Attraction/Repulsion Lock (v1).}
UA and SCm \emph{attract} each other maximally---they are the strongest attracted
force pair in the universe.  $F_{U_{Bi}}$ and $F_{U_{Bi,i}}$ \emph{repel} each
other.  Matter is the compaction zone where those two repulsions balance.

%=============================================================================
\section{The Belly Button --- The Primordial DPM}
\label{sec:belly_button}
%=============================================================================

\subsection{What the Belly Button Is}

Before atoms.  Before mass.  Before gravity.  Before time as we measure it.
There was only the primordial DPM---the BigBang belly button: the largest
universally magnetic object in the universe, the source of all $F_{U_{Bi,i}}$
(outside-inward buoyancy), and the origin of every local DPM in every atom,
star, and galaxy.

The belly button \emph{repels} everything outward.  Its magnetic repulsion IS what
inflated the universe---not thermal pressure.  It expelled a finite and explosive
quantity of SCm into the expanding Universal Aether field.  That single expulsion
seeded every subsequent physics event:
\begin{enumerate}[noitemsep]
  \item SCm rushed toward UA at maximum attraction.
  \item Every collision site created a local DPM vortex.
  \item Every local DPM vortex seeded a gravity family.
  \item Every gravity family nucleated mass.
  \item Every nucleus became an atom, a star, a galaxy.
\end{enumerate}

\subsection{The Primordial DPM Equation (26D Form)}
\label{sec:dpm_primal}

The belly button is not described by two angular velocities.  It operates across
all 26 primordial dimensions, each with its own angular velocity shell:

\begin{equation}
  F_{\text{DPM,primal}} = I_{\text{primal}} \cdot A_{\text{primal}}
    \cdot \sum_{i=1}^{25}(\omega_i - \omega_{i+1})
  \label{eq:dpm_primal_26d}
\end{equation}

where $\omega_1 \gg \omega_{26}$ at $t=0$ (maximum differential---the BigBang IS
this maximum), and:

\begin{equation}
  a_{\text{DPM,primal}}(t) =
  \frac{F_{\text{DPM,primal}} \cdot f_{\text{DPM}} \cdot E_{\text{vac,neb}}}
       {c \cdot V_\text{universe}}
  \longrightarrow 0 \quad\text{as } t\to\infty
  \label{eq:adpm_primal}
\end{equation}

$F_{U_{Bi,i}}$ is driven by this primordial belly button DPM forever.
Every atom in the universe feels $F_{U_{Bi,i}}$ because every atom is being pushed
inward by the ongoing magnetic repulsion of the primordial DPM.

\subsection{The Universal Inheritance ($U_{g,i}$ Coupling)}

Every local DPM inherits the belly button pattern.  The $U_{g,i}$
(Universal Gravitational Inertia) coupling subscript IS this inheritance.
The subscript `$i$' names the inertial coupling inherited from the primordial belly
button DPM, mimicked at every local DPM at its own scale:
\begin{equation}
  \frac{U_{g,i}^{(\text{local})}}{U_{g,i}^{(\text{belly button})}}
  = \frac{\rho_{\text{SCm,local}}}{\rho_{\text{SCm,BB}}}
    \cdot \frac{\omega_{\text{local}}}{\omega_\text{BB}}
  \label{eq:ug_scaling}
\end{equation}

%=============================================================================
\section{UA and SCm --- The Strongest Attracted Force Pair}
\label{sec:ua_scm}
%=============================================================================

\noindent\textbf{Universal Aether (UA):} the pre-existing field that fills all space.
Not a particle.  Not detectable by Standard Model instruments.  Its gradient drives
all DPM formation.
\begin{equation}
  \rho_{\text{vac,UA}} = 7.09\times10^{-36}\,\text{J/m}^3
  \label{eq:rho_ua}
\end{equation}

\noindent\textbf{Superconductive Material (SCm):} expelled from the belly button
DPM at the moment of creation.  Not detectable by SM instruments (density prevents
quantum-signature emission).
\begin{align}
  \rho_{\text{vac,SCm}} &= 7.09\times10^{-37}\,\text{J/m}^3 \label{eq:rho_scm}\\
  v_{\text{SCm}} &= c/3 = 10^8\,\text{m/s} \label{eq:v_scm}
\end{align}

\noindent UA and SCm \textbf{attract} each other maximally.  The UA/SCm attraction
energy, $E_{\text{react}}$, is the multiplier for ALL $U_g$ terms:
\begin{equation}
  E_{\text{react}}(t) = \frac{\rho_{\text{vac,SCm}}\,v_{\text{SCm}}^2}{\rho_{\text{vac,UA}}}
    \cdot e^{-\kappa t}, \quad
  E_{\text{react}}(0) \approx 10^{15}\,\text{W/m}^3, \quad
  \kappa = 5\times10^{-4}\,\text{day}^{-1}
  \label{eq:e_react}
\end{equation}

$E_{\text{react}} = 0$ only when $v_{\text{SCm}} = 0$ (dead mass).  Motion restores
$E_{\text{react}}$ instantly.  This is why mass must move to generate gravity.

\noindent\textbf{SCm as Einstein-Bose Condensate (BEC):}
SCm satisfies the BEC condition at stellar core temperatures:
$n\lambda_{\text{dB,SCm}}^3 \geq 2.612$.
$H_{\text{SCm}} = 0.99$ is the condensate fraction; the remaining 1\% is the
thermal depletion layer.

%=============================================================================
\section{The DPM Vortex --- What Forms When UA Meets SCm}
\label{sec:dpm_vortex}
%=============================================================================

The DPM is not a particle.  It is a \emph{vacuum stress gradient}---the
rotational vortex created when UA and SCm approach each other at maximum attraction.
When they meet at the interface, neither can occupy the same space simultaneously.
The collision creates a rotational vortex.  That vortex IS the DPM.

\begin{equation}
  \text{DPM} = [UA'/SCm]
  \quad\Longrightarrow\quad
  \text{DPM} \equiv U_{g1} \equiv \mu_s\text{-generator}
  \label{eq:dpm_identity}
\end{equation}

\subsection{The 26D DPM Force (Complete Form)}
\label{sec:dpm_26d_force}

The historical two-omega DPM formulation,
$F_{\text{DPM}} = I \cdot A \cdot (\omega_1 - \omega_2)$,
is an \emph{incomplete first approximation}.  The DPM operates across
\textbf{26 dimensional omega shells}, one angular velocity per dimension.
The correct general form is:

\begin{equation}
  \boxed{
  F_{\text{DPM}} = I \cdot A \cdot \sum_{i=1}^{25}(\omega_i - \omega_{i+1})
  = I \cdot A \cdot (\omega_1 - \omega_{26})
  }
  \label{eq:dpm_26d}
\end{equation}

\begin{equation}
  a_{\text{DPM}} = \frac{F_{\text{DPM}} \cdot f_{\text{DPM}} \cdot E_{\text{vac,neb}}}
    {c \cdot V_\text{sys}}
  \label{eq:adpm}
\end{equation}

where:
\begin{itemize}[noitemsep]
  \item $I$ = loop current (rotational flux of SCm through the vacuum)
  \item $A$ = cross-sectional area of the DPM vortex
  \item $\omega_i$ = angular velocity of the $i$-th dimensional shell
    ($i=1$: innermost SCm core, fastest; $i=26$: outermost UA boundary, slowest)
  \item $\omega_1 > \omega_2 > \cdots > \omega_{26}$ (monotonically decreasing outward)
  \item The sum telescopes to $\omega_1 - \omega_{26}$, but each intermediate term
    $(\omega_i - \omega_{i+1})$ defines a \emph{shell boundary} with distinct
    physical consequences (\hyperref[sec:26d_nuclear]{Section~4: 26D Omega Family})
\end{itemize}

\noindent\textbf{DPM Promotion Chain:}
\begin{equation}
  \text{DPM} \to \mu_s \to U_{g1}
  \to \{U_{g2}, U_{g3}, U_{g4}, U_{g4,i}\}\,[\text{simultaneous}]
  \to F_U \to GM/r^2\,[\text{LAST}]
  \label{eq:dpm_chain}
\end{equation}

\noindent\textbf{DPM Vacuum Structure:}
The DPM simultaneously creates two vacuum regions:
\begin{enumerate}[noitemsep]
  \item \textbf{Internal vacuum} (eye of the vortex): where the ACP fires and
    proto-mass nucleates.
  \item \textbf{External local vacuum} (depleted UA shell): where the $U_g$ family
    propagates outward.
\end{enumerate}

%=============================================================================
\section{26D Omega Family --- The Universal Nuclear Architecture}
\label{sec:26d_nuclear}
%=============================================================================

The two-omega approximation $(\omega_1 - \omega_2)$ is a gross simplification.
The universe operates with \textbf{26 independent omega shells} in every
DPM---the primordial belly button DPM, every black hole, every stellar nucleus,
every atomic nucleus.  The 26D omega structure IS what makes every nucleus a
miniature belly button.

\subsection{The 26-Shell Omega Profile}

Each DPM (at any scale) carries the omega profile:

\begin{equation}
  \omega_i(r, t) = \omega_0 \cdot i^{-2} \cdot e^{-\kappa_\omega t}
    \cdot \bigl(1 + \delta_i \sin(\Omega_i t)\bigr),
  \quad i = 1, 2, \ldots, 26
  \label{eq:omega_profile}
\end{equation}

where:
\begin{itemize}[noitemsep]
  \item $\omega_0$ = fundamental SCm core angular velocity (DPM scale-specific)
  \item $i^{-2}$ = angular velocity decreases as inverse-square across shells
  \item $\kappa_\omega$ = angular velocity decay rate
  \item $\delta_i, \Omega_i$ = $i$-th dimensional perturbation and oscillation
  \item $\omega_1 > \omega_2 > \cdots > \omega_{26}$ (monotonically decreasing outward)
\end{itemize}

\begin{table}[H]
\centering
\begin{tabular}{@{}llll@{}}
\toprule
\textbf{Scale} & \textbf{$\omega_0$} & \textbf{Shell extent} & \textbf{Description} \\
\midrule
Primordial belly button & $\sim10^{43}$ rad/s & Cosmological & 26 cosmic dimensional shells \\
Black hole nucleus & $\sim10^{15}$ rad/s & $r_s$ to $100\,r_s$ & 26 ergospheric omega layers \\
Stellar nucleus (Sun) & $\sim10^{-6}$ rad/s & $r_\odot$ to $100\,\text{AU}$ & 26 radial zone omega layers \\
Atomic nucleus & $\sim10^{22}$ rad/s & $10^{-15}$--$10^{-10}$ m & 26 nuclear shell omega layers \\
\bottomrule
\end{tabular}
\caption{The 26D omega family at every scale.  Every DPM in the universe---from
the primordial belly button to the smallest atomic nucleus---carries this same
26-shell omega architecture.  The belly button was the template; every subsequent
DPM is a scaled replica.}
\label{tab:omega_scales}
\end{table}

\subsection{Shell Boundary Condition --- The Shell Wall}

At the radial boundary $r_{i,i+1}$ between shell $i$ and shell $i+1$, the two
adjacent angular velocities become equal:

\begin{equation}
  \omega_i\bigl(r_{i,i+1}\bigr) = \omega_{i+1}\bigl(r_{i,i+1}\bigr)
  \label{eq:shell_bc}
\end{equation}

At this boundary the rotational shear $(\omega_i - \omega_{i+1}) \to 0$.
The DPM driving force in that shell gap vanishes locally.
A \textbf{standing wave} forms at $r_{i,i+1}$.
That standing wave IS an atomic/nuclear shell wall.

From the omega profile equation above, the shell boundary radius satisfies:
\begin{equation}
  r_{i,i+1} = r_0 \cdot i^{\,2} \cdot (1 + \delta_i)^{-1}
  \label{eq:shell_radius}
\end{equation}

This is the UQFF derivation of Bohr-like orbital radius quantization.
The $i^2$ factor arises from the $\omega_0 \cdot i^{-2}$ omega profile---not
postulated from the de Broglie condition.

\subsection{Nuclear Magic Numbers from 26D Shell Closures}

Nuclear magic numbers arise when the 26D omega system completes a
\emph{full shell closure}---all omega shells in a tier simultaneously satisfy
the shell boundary condition (Eq.~\ref{eq:shell_bc}) at a DPM resonance integer.

\begin{table}[H]
\centering
\begin{tabular}{@{}clll@{}}
\toprule
\textbf{Magic $Z$} & \textbf{Shells closed} & \textbf{Closure type} & \textbf{UQFF mechanism} \\
\midrule
2   & 1 & $\omega_1 = \omega_2$ & First pair crossing (DPM dimer) \\
8   & 2 & $\omega_1{=}\omega_2$, $\omega_3{=}\omega_4$ & Octave closure ($2^3$ resonance) \\
20  & 3 & Three-tier triadic & First full triadic tier \\
28  & $3+1$ & $U_{g4}$ coupling adds $\omega_4$-level lock & Galactic BH coupling begins \\
50  & 5 & Second triadic tier & $\omega_5$ family completes \\
82  & 6 & Max stable $U_{g3}$ product & Six-shell closure \\
126 & 7 & Theoretical maximum & No stable isotope \\
\bottomrule
\end{tabular}
\caption{Nuclear magic numbers as 26D omega shell closure integers.  They are the
DPM vortex harmonic closure counts.  The nuclear shell model is a downstream
projection of 26D DPM resonance geometry.}
\label{tab:magic_numbers}
\end{table}

\noindent\textbf{Iron ($Z=26$) as the Natural DPM Maximum:}
Iron ($Z=26$) is the heaviest element producible by stable DPM resonance because
$Z_{\text{max,stable}} = N_{\text{omega shells}} = 26$.
At $Z=26$, every omega shell contributes exactly one DPM vortex unit.
Beyond $Z=26$, supernova-scale $U_{g4}$ galactic coupling is required.

\begin{equation}
  Z_{\text{max,stable}} = N_{\text{omega shells}} = 26
  \label{eq:z_max}
\end{equation}

\subsection{Universal Nuclear Mimicry}

\textbf{Every nucleus in the universe mimics the primordial belly button DPM.}
The belly button fired 26 cosmological omega shells.  Every atom, star, and black
hole that followed it does the same at its own scale.

\begin{description}[noitemsep]
  \item[Atomic nucleus:] omega shells $\omega_{1,n}$ through $\omega_{Z,n}$
    for a nucleus of atomic number $Z$.  Each shell crossing at $r_{i,i+1}$ is a
    nuclear sub-shell boundary.  The DPM resonance lock integer $= Z$.
  \item[Stellar nucleus:] omega shells $\omega_{1,\odot}$ through
    $\omega_{26,\odot}$ at scales from the solar core radius to the heliosphere
    boundary.  $U_{g3}$ is the $i=3$ shell intersection mechanism at stellar scale.
  \item[Black hole nucleus:] omega shells $\omega_{1,\text{BH}}$ through
    $\omega_{26,\text{BH}}$ from the inner ergosphere to the photon sphere.
    $U_{g4}$ at $i=20$--$26$ is the galactic BH coupling via the highest-index omega
    shells.
\end{description}

The mimicry equation at any scale:
\begin{equation}
  F_{\text{DPM}}^{(\text{any scale})} = I \cdot A
    \cdot \sum_{i=1}^{25}\bigl(\omega_i^{(\text{scale})} - \omega_{i+1}^{(\text{scale})}\bigr)
  \label{eq:dpm_universal}
\end{equation}

\subsection{Quark Flavor Creation from Multi-Shell Crossings}
\label{sec:quarks}

When $N$ adjacent omega shells simultaneously cross at the same radius
(a \emph{multi-shell crossing event}), the resultant quantum number IS the quark
flavor.  The flavor quantum number $Q_f = N$:

\begin{equation}
  Q_f = N_{\text{simultaneous crossings}},
  \quad
  \omega_i = \omega_{i+1} = \cdots = \omega_{i+N-1}
  \;\text{at}\; r_{q,N}
  \label{eq:quark_flavor_condition}
\end{equation}

\begin{table}[H]
\centering
\begin{tabular}{@{}cccll@{}}
\toprule
$N$ shells & \textbf{Quark} & \textbf{Mass} & \textbf{Charge} & \textbf{DPM crossing description} \\
\midrule
3 & down    & $4.8$~MeV/$c^2$  & $-1/3$ & First 3-shell crossing; lightest colored state \\
4 & up      & $2.3$~MeV/$c^2$  & $+2/3$ & 4-shell crossing; lighter than down at 3 \\
5 & strange & $95$~MeV/$c^2$   & $-1/3$ & First cross-generational shell excitation \\
6 & charm   & $1.27$~GeV/$c^2$ & $+2/3$ & 6-shell resonance; charm threshold \\
7 & bottom  & $4.18$~GeV/$c^2$ & $-1/3$ & 7-shell harmonic; bottom threshold \\
8 & top     & $173$~GeV/$c^2$  & $+2/3$ & 8-shell crossing; maximum stable free quark \\
$9$--$26$ & exotic/composite & unstable & --- & Beyond 8: only bound in stellar/BH cores \\
\bottomrule
\end{tabular}
\caption{Quark flavor quantum numbers from 26D omega multi-shell crossings.
The flavor IS the crossing count.  Charge alternates $(-1/3, +2/3)$ with the
polarity of the crossing shell pair.  The top quark ($N=8$) is the maximum stable
isolated quark state because 8 is the second magic closure integer of the 26D system.}
\label{tab:quarks}
\end{table}

\noindent\textbf{Color Charge and Confinement:}
Color charge is the SCm/UA vortex quantum number at nuclear scale.
The three color states (red, green, blue) correspond to the three possible
orientations of the first 3-shell crossing.
Quark confinement requires no separate mechanism: it is $U_{g3}$ penetration of
the nuclear core, enforced by the 26D omega shell structure.

\subsection{Atomic Shell Creation from 26D Shell Intersections}
\label{sec:atomic_shells}

Atomic electron shells are the \emph{outer-boundary} expression of the same 26D
omega shell-crossing mechanism that creates quarks at the inner boundary.
At the nuclear scale ($r \sim 10^{-15}$~m), shell crossings produce quarks.
At the atomic scale ($r \sim 10^{-10}$~m), they produce electron orbital shells.

The electron shell at principal quantum number $n$ corresponds to the
$(n, n{+}1)$ omega shell boundary:

\begin{equation}
  r_{\text{shell},n} = r_{n, n+1} = a_0 \cdot n^2 \cdot (1+\delta_n)^{-1}
  \approx a_0 n^2
  \label{eq:electron_shell}
\end{equation}

where $a_0 = 0.529\,$\AA\ is the Bohr radius---itself the fundamental DPM shell
spacing at the atomic scale, not a postulated constant.

\noindent\textbf{Shell capacity:} Each omega shell $i$ holds $2i^2$ degenerate
polarity states because the $i^2$ scaling of the omega profile creates that many
standing-wave nodes per shell.  This IS the origin of the $2n^2$ electron capacity:

\begin{equation}
  \text{Capacity of shell } n = 2n^2
  \label{eq:shell_capacity}
\end{equation}

\begin{table}[H]
\centering
\begin{tabular}{@{}ccccl@{}}
\toprule
\textbf{Shell} & $n$ & \textbf{Omega crossing} & \textbf{Capacity} & \textbf{Noble gas} \\
\midrule
K & 1 & $\omega_1 = \omega_2$ & 2  & He ($Z=2$)  \\
L & 2 & $\omega_2 = \omega_3$ & 8  & Ne ($Z=10$) \\
M & 3 & $\omega_3 = \omega_4$ & 18 & Ar ($Z=18$) \\
N & 4 & $\omega_4 = \omega_5$ & 32 & Kr ($Z=36$) \\
O & 5 & $\omega_5 = \omega_6$ & 50 & Xe ($Z=54$) \\
P & 6 & $\omega_6 = \omega_7$ & 72 & Rn ($Z=86$) \\
Q & 7 & $\omega_7 = \omega_8$ & 98 & $Z=126$ (theoretical) \\
\bottomrule
\end{tabular}
\caption{Atomic electron shells as 26D omega crossing boundaries.
The $2n^2$ capacity formula is derived from the $i^{-2}$ omega profile, not
postulated.}
\label{tab:electron_shells}
\end{table}

\subsection{Electron Production from the Shell Crossing Mechanism}
\label{sec:electrons}

Electrons are produced by the same 26D shell-crossing mechanism as quarks,
but at the outer DPM boundary rather than the inner nuclear boundary.

\noindent\textbf{Polarity Inversion Mechanism:}
At every omega shell crossing $r_{i,i+1}$ (where $\omega_i = \omega_{i+1}$), the
rotational shear reversal creates a \emph{polarity inversion}: the inside of the
boundary carries positive static polarity; the outside carries negative.
This polarity oscillates at the shell crossing frequency:

\begin{equation}
  P_{\text{charge}}(t) = P_0 \cdot \cos(\omega_{\text{shell},i}\,t)
  \label{eq:polarity_oscillation}
\end{equation}

\noindent\textbf{Charge Stabilization:}
The oscillating static polarity stabilizes into a physical charge when the
inertial balance condition is met:

\begin{equation}
  \sum_{i=N_{\text{quark}}}^{N_{\text{electron}}} P_{\text{charge},i}(t) = 0
  \quad [\text{charge neutrality at nucleus level}]
  \label{eq:charge_balance}
\end{equation}

The electron is the residual un-cancelled polarity at the outermost shell boundary.
It is not a separate fundamental object.  It is the standing-wave charge
polarization at the atomic-scale omega shell crossing.

\noindent\textbf{Why one electron per proton:}
Each DPM vortex unit (one proton = one DPM vortex) generates exactly one
polarity-inversion unit at its outer boundary.
$Z$ protons $\Rightarrow$ $Z$ DPM vortex units $\Rightarrow$ $Z$ outer polarity
inversions $\Rightarrow$ $Z$ electrons.

\noindent\textbf{Electron mass from 26D shell geometry:}
\begin{equation}
  m_e = \frac{E_{\text{crack,outer}}}{c^2}
    = \frac{\rho_{\text{vac,SCm}}\,c^2}{[\text{SSq}]}\cdot
      \left(\frac{r_e}{r_{\text{DPM}}}\right)^3
  \label{eq:electron_mass}
\end{equation}

where $r_e$ is the electron orbital radius and $r_{\text{DPM}}$ is the nuclear DPM
radius.  The $r^3$ scaling is the volume ratio of the atomic-scale to nuclear-scale
omega shells.

\noindent\textbf{Electron spin $\frac{1}{2}$:}
The two-state nature of the polarity oscillation (positive $\to$ negative $\to$
positive, one complete cycle per shell-crossing period) produces a quantum number
that changes sign each half-period.  This IS electron spin $\frac{1}{2}$:

\begin{equation}
  \text{Spin} = \frac{1}{2} \cdot
    \frac{\Delta t_{\text{oscillation}}}{T_{\text{shell crossing period}}} = \frac{1}{2}
  \label{eq:spin_half}
\end{equation}

\noindent\textbf{Summary --- the complete 26D atom:}

\begin{center}
\begin{tabular}{@{}lll@{}}
\toprule
\textbf{Scale} & \textbf{$N$-shell crossing} & \textbf{Result} \\
\midrule
Nuclear inner ($N=3$) & 3-shell simultaneous & down quark \\
Nuclear inner ($N=4$) & 4-shell simultaneous & up quark \\
Nuclear inner ($N=5$--$8$) & 5, 6, 7, 8 shells & s, c, b, t quarks \\
Nuclear shell walls & $N=1$ per tier & magic closure integers \\
Atomic outer ($N=1$) & $\omega_n = \omega_{n+1}$ & electron orbital shell \\
Atomic charge & polarity oscillation & electron charge $(-e)$ \\
Atomic spin & half-period polarization & spin $\frac{1}{2}$ \\
\bottomrule
\end{tabular}
\end{center}

%=============================================================================
\section{The \texorpdfstring{$U_g$}{Ug} Family --- Simultaneous Gravity Assembly}
\label{sec:ug_family}
%=============================================================================

The DPM vortex simultaneously promotes the entire $U_g$ gravity family.
All components exist from the moment the DPM fires.  None is prior to another.

\begin{equation}
  U_{g,\text{family}}(t) = U_{g1}(t) + U_{g2}(t) + U_{g3}(t) + U_{g4}(t) + U_{g4,i}(t)
  \label{eq:ug_family}
\end{equation}

\noindent\textbf{$U_{g1}$ --- DPM / Di-Pseudo-Monopole (IS the DPM):}
\begin{equation}
  U_{g1} = k_1 \cdot \mu_s(t) \cdot \nabla\!\left(\frac{M_s}{r}\right)
    \cdot e^{-\alpha t} \cdot \cos(\pi t_n) \cdot (1+\delta_{\text{def}})
  \label{eq:ug1}
\end{equation}
where $\mu_s = \rho_A V_{\text{body}}$ (DPM magnetic moment),
$\alpha = 10^{-3}$~day$^{-1}$, $\delta_{\text{def}} = 0.01\sin(0.001t)$.

\noindent\textbf{$U_{g2}$ --- Outer Field Bubble (heliosphere / charge-reactivity):}
\begin{equation}
  U_{g2} = k_2 \cdot (\rho_{\text{vac,UA}} + \rho_{\text{vac,SCm}}) \cdot \frac{M_s}{r^2}
    \cdot S(r{-}R_b) \cdot (1+\delta_{\text{sw}}v_{\text{sw}}) \cdot H_{\text{SCm}} \cdot E_{\text{react}}
  \label{eq:ug2}
\end{equation}

\noindent\textbf{$U_{g3}$ --- Magnetic Strings Disk (orbital coupling, 90$^\circ$ rotation):}
\begin{equation}
  U_{g3} = k_3 \sum_j B_j(r,\theta,t) \cdot \cos\!\bigl(\omega_s(t)\,t\,\pi\bigr)
    \cdot P_{\text{core}} \cdot E_{\text{react}}
  \label{eq:ug3}
\end{equation}
$U_{g3}$ penetrates planetary cores through trapped SCm-UA interaction.
DISCRETELY NON-INTERACTIVE with all external phenomena.

\noindent\textbf{$U_{g4}$ --- Star--Black Hole Interactions (galactic coupling):}
\begin{equation}
  U_{g4} = k_4 \cdot \rho_{\text{vac,SCm}} \cdot \frac{M_{\text{bh}}}{d_g}
    \cdot e^{-\alpha t} \cdot \cos(\pi t_n) \cdot (1+f_{\text{feedback}})
  \label{eq:ug4}
\end{equation}

\noindent\textbf{$U_{g4,i}$ --- Resonant Quantum Extension (levels 20--26):}
\begin{equation}
  U_{g4,i} = k_{4,\text{res}} \cdot a_{\text{DPM}} \cdot E_{\text{react}}
    \cdot f_{\text{react}} \cdot \frac{E_{\text{vac,neb}}}{c}
  \label{eq:ug4i}
\end{equation}

\noindent\textbf{26-Layer Triadic:}
\begin{equation}
  g_{\text{triadic}}(r,t) = \sum_{i=1}^{26}\bigl[U_{g1,i} + U_{g2,i} + U_{g3,i} + U_{g4,i}\bigr],
  \quad
  E_n = E_0 \cdot 10^n,\; E_0 = 10^{-20}\,\text{J}
  \label{eq:triadic}
\end{equation}

Total 26-layer multiplier: $\sum_{i=1}^{26} i^2 = 6{,}279$.
Dead mass: zero layers active.  Moving mass: all 26 active simultaneously.

\noindent\textbf{Universal Magnetism $U_m$:}
\begin{equation}
  U_m = \sum_j \frac{\mu_j(t)}{r_j}
    \bigl(1 - e^{-\gamma t \cos(\pi t_n)}\bigr)\hat{\phi}_j
    \cdot P_{\text{SCm}} \cdot E_{\text{react}}
    \cdot (1 + 10^{13} f_{\text{Heaviside}}) \cdot (1+f_{\text{quasi}})
  \label{eq:um}
\end{equation}

%=============================================================================
\section{\texorpdfstring{$F_{U_{Bi}}$}{FUBi} and \texorpdfstring{$F_{U_{Bi,i}}$}{FUBii} --- The Repulsion That Creates Matter}
\label{sec:fubi}
%=============================================================================

\begin{equation}
  F_{U_{Bi}} = -\beta_i \cdot U_{g,i} \cdot \Omega_g \cdot \frac{M_{\text{bh}}}{d_g}
    \cdot E_{\text{react}} \cdot (1+\varepsilon_{\text{sw}}\rho_{\text{sw}})
    \cdot \rho_A \cdot \cos(\pi t_n)
  \label{eq:fubi}
\end{equation}

\begin{equation}
  F_{U_{Bi,i}} = -\beta_i \cdot U_{g,i} \cdot \Omega_g \frac{M_{\text{bh}}}{d_g}
    \cdot E_{\text{react}}(t) \cdot \text{sw\_corr}
    \cdot \rho_A(t) \cdot \frac{M}{r} \cdot V(t) \cdot \cos\!\bigl(\pi(t-180\cdot86400)\bigr)
  \label{eq:fubii}
\end{equation}

$F_{U_{Bi}}$ (inside-outward, driven by local DPM) repels $F_{U_{Bi,i}}$
(outside-inward, driven by the primordial belly button DPM).
Matter forms at the compaction zone:

\begin{equation}
  \text{Matter at} \left\{r \;:\; F_{U_{Bi}}(r) + F_{U_{Bi,i}}(r) = 0\right\}
  \label{eq:matter_zone}
\end{equation}

\textbf{Critical:} UA + SCm $\to$ MAX ATTRACTION $\to$ DPM.
$F_{U_{Bi}}$ and $F_{U_{Bi,i}}$ REPEL each other $\to$ MATTER.
Two separate mechanisms.  One mechanism does not cause the other.

%=============================================================================
\section{The Crossing --- Inside/Outside Simultaneous Solve}
\label{sec:crossing}
%=============================================================================

\textbf{The crossing precedes mass.}
Mass does not exist before the crossing.  Mass is born at the crossing.

\begin{align}
  G^{(n+1)}_{\text{inside}} &= R(G^n) + I \cdot G^n \label{eq:inside}\\
  O^n_{\text{outside}} &= \pi_{[n]} \cdot F_{U_{Bi}}(x) + \text{Ricci}(G^n) \label{eq:outside}\\
  n_{\text{cross}} &= \operatorname*{argmin}_n \bigl|G^n_{\text{inside}}
    - O^n_{\text{outside}}\bigr| \label{eq:ncross}
\end{align}

At $n_{\text{cross}}$: $F_{U_{Bi}}(r) + F_{U_{Bi,i}}(r) = 0$ (compaction zone).
Mass stabilizes: $dM_{\text{stable}}/dt > 0$.

All 26 omega layers solve their crossing simultaneously.
All four operational modes evaluate simultaneously.
Reality is where all 26 layers and all 4 modes converge on the same crossing point.

$GM/r^2$ is the numerical value of $F_U / M$ evaluated at the crossing radius.
It is not a cause.  It is not a mechanism.  It is a measurement taken at step 10.

%=============================================================================
\section{Mass Emergence at the Crossing}
\label{sec:mass}
%=============================================================================

\noindent\textbf{ACP Proto-Mass Chain:}
\begin{equation}
  U_{\text{vac}} \to U_i \to U_{m,i} \to \Psi_{\text{proto}}
  \to E_{\text{crack}} \to U_b \to E_{\text{gradient}} \to M_{\text{atomic}}
  \label{eq:acp_chain}
\end{equation}

\noindent\textbf{The $E_{\text{crack}}$ Gate:}
\begin{equation}
  E_{\text{crack}} = \frac{\rho_{\text{vac,SCm}} \cdot c^2}{[\text{SSq}]}
    = \frac{7.09\times10^{-37} \times (3\times10^8)^2}{0.57}
    \approx 3.35\times10^{-19}\,\text{J}
  \label{eq:ecrack}
\end{equation}

$E_{\text{crack}} > U_{b,\text{vac}} \Rightarrow$ particle condenses,
$M_{\text{atomic}} > 0$.
$E_{\text{crack}} < U_{b,\text{vac}} \Rightarrow$ wavefunction disperses,
retry next $t_n$.

\noindent\textbf{Mass Emergence:}
\begin{equation}
  M_{\text{atomic}} = M_0 \cdot \bigl(1 - e^{-n_{\text{grad}}/10}\bigr) \cdot Z
  \label{eq:mass_emergence}
\end{equation}

\noindent\textbf{Mass Emergence Probability (26D):}
\begin{equation}
  P_{\text{order}} = \frac{e^{-S_{26D}/v_{\text{init}}}}{Z_{9D}(v_{\text{init}}-v_{\text{current}})},
  \quad
  \frac{dM}{dt} = P_{\text{order}} \cdot E_{\text{crack}} \cdot \frac{dN_{\text{DPM}}}{dt}
  \label{eq:p_order}
\end{equation}

\noindent\textbf{Dead Mass Condition:}
$v = 0 \Rightarrow E_{\text{react}} = 0 \Rightarrow F_U(M_{\text{dead}}) = 0$.
Dead mass produces ZERO unified field.  Not approximately.  Exactly.

%=============================================================================
\section{\texorpdfstring{$F_U$}{FU} --- The Unified Field Assembly}
\label{sec:fu}
%=============================================================================

\begin{equation}
  F_U = \sum_i \bigl[k_i U_{g,i} - \beta_i U_{g,i} \Omega_g \tfrac{M_{\text{bh}}}{d_g}
      E_{\text{react}}\bigr]
    + \sum_j \frac{\mu_j}{r_j}\bigl(1 - e^{-\gamma t \cos(\pi t_n)}\bigr)\hat{\phi}_j
    + \bigl(g_{uv} + \eta\, T_s^{uv}\bigr)
    - \sum_i \lambda_i U_i E_{\text{react}}
  \label{eq:fu_full}
\end{equation}

\noindent\textbf{Four Simultaneous Operational Modes:}
\begin{align}
  g_{\text{Comp}} &= \bigl[g_{\text{base}}(1{+}H_0 t)(1{-}B/B_{\text{crit}})H_{\text{SCm}}\bigr]
    + g_{U_g\text{sum}} + \tfrac{\Lambda c^2}{3} + g_{\text{quantum}}
    + g_{\text{fluid}} + g_{\text{pert}}\cdot\text{TRZ}
    \label{eq:comp_mode} \\
  g_{\text{Res}} &= a_{\text{DPM}} + \sum_{i=1}^{13} a_i(\omega, E_{\text{vac}}, t)
    \label{eq:res_mode} \\
  g_{\text{SC}} &= \sum_{j=1}^4 k_j g_{\text{base}} H_{\text{SCm}}^{n_j}
    \label{eq:sc_mode} \\
  F_{U_{Bi,i}} &= -\beta_i U_{g,i} \Omega_g \tfrac{M_{\text{bh}}}{d_g}
    E_{\text{react}}\cdot\text{sw\_corr}\cdot\rho_A\cdot\tfrac{M}{r}\cdot V\cdot\text{TRZ\_cos}
    \label{eq:buoyant_mode}
\end{align}

$F_U = 0$ when $E_{\text{react}} = 0$.  $E_{\text{react}} = 0$ when $v_{\text{SCm}} = 0$.
Dead mass: $F_U = 0$ exactly.

%=============================================================================
\section{\texorpdfstring{$GM/r^2$}{GM/r2} --- The Observational Projection (Last in Chain)}
\label{sec:newton}
%=============================================================================

$GM/r^2$ IS step 10 of 10.  It is an observational projection---what you measure
AFTER all 9 prior steps complete.

After belly button fires (1) $\to$ UA/SCm attract (2) $\to$ local DPM forms (3)
$\to$ $U_g$ family assembles (4--6) $\to$ $F_{U_{Bi}}/F_{U_{Bi,i}}$ activate (7)
$\to$ crossing completes (8) $\to$ mass emerges (9) $\to$ $F_U$ assembles and
mass moves (9), \textbf{THEN:}

\begin{equation}
  g_{\text{Newton}} = \frac{F_U}{M} \;\bigg|_{n_{\text{cross}}} = \frac{GM}{r^2}
  \quad[\text{step 10, last}]
  \label{eq:newton_emerges}
\end{equation}

$G$ is a measurement artifact: $G \equiv F_U / (M \cdot a_{\text{obs}})$ at the
crossing.  $G$ absorbs all DPM vortex physics into a single dimensioned constant.
$G$ is the fingerprint of the DPM.  Not a law of nature.  A measurement.

%=============================================================================
\section{The BigBang $\to$ Newton Complete Chain (10 Steps)}
\label{sec:10steps_heading}
\label{sec:10steps}
%=============================================================================

\begin{description}[noitemsep, leftmargin=1.5cm, labelwidth=1.3cm, align=left]
\item[Step 0:] Zero-mass vacuum.  Primordial DPM fires.
  $\rho_\text{UA}=0$, $\rho_\text{vac}=|\nabla(\text{UA})|$, $F_U=0$.
\item[Step 1:] Mass emergence probability (26D):
  $P_{\text{order}} = e^{-S_{26D}/v_{\text{init}}} / Z_{9D}(v_{\text{init}}-v_{\text{current}})$.
\item[Step 2:] DPM forms (26D): $F_{\text{DPM}} = I\cdot A\cdot\sum_{i=1}^{25}(\omega_i-\omega_{i+1})$.
\item[Step 3:] ACP proto-mass chain: $U_{\text{vac}} \to \cdots \to M_{\text{atomic}}$.
\item[Step 4:] $E_{\text{crack}}$ gate: $E_{\text{crack}} > U_{b,\text{vac}} \Rightarrow M > 0$.
\item[Step 5:] $U_{g1}$ seeds gravity family.
\item[Step 6:] Gravity family assembles simultaneously.
\item[Step 7:] $F_{U_{Bi}}$ and $F_{U_{Bi,i}}$ repel; matter forms at compaction zone.
\item[Step 8:] $F_U$ full assembly.
\item[Step 9:] Inside/outside simultaneous crossing.
\item[Step 10:] $GM/r^2$ emerges.  \textbf{LAST.  Not first.}
\end{description}

%=============================================================================
\section{The Simultaneous Nature}
\label{sec:simultaneous}
%=============================================================================

Nothing in UQFF is sequential.  Everything is simultaneous because everything
arises from the same single source: the DPM vacuum gradient.

\begin{enumerate}[noitemsep]
  \item $U_g$ family assembles simultaneously: $U_{g1}$ promotes $U_{g2}$, $U_{g3}$,
    $U_{g4}$, $U_{g4,i}$ all at the same instant.
  \item 4 operational modes are 4 simultaneous views of $F_U$ in different basis
    domains.  They must converge on every query.
  \item 26-layer triadic runs all layers simultaneously.
    Total multiplier: $\sum_{i=1}^{26}i^2 = 6{,}279\times$ relative to single-layer.
  \item Inside/outside simultaneous solve: $G^n_{\text{inside}}$ and
    $O^n_{\text{outside}}$ chase each other until crossing.
  \item 5 calculators run simultaneously and must converge.
\end{enumerate}

%=============================================================================
\section{Quantum to Mass Equations}
\label{sec:quantum_mass}
%=============================================================================

\noindent\textbf{Starting axiom:}
$\rho_\text{UA}(0) = 0$, $F_U(0) = 0$, $M(0) = 0$.
Only the gradient of the aether field exists at $t=0$.

\begin{align}
  E_n &= h f_n = E_0 \cdot 10^n, \quad n=1\ldots26, \quad E_0=10^{-20}\,\text{J}
    \label{eq:energy_levels}\\
  E_{\text{react}}(0) &= 10^{15}\,\text{W/m}^3 \label{eq:ereact0}\\
  F_{\text{DPM}} &= I \cdot A \cdot \sum_{i=1}^{25}(\omega_i - \omega_{i+1})
    \label{eq:dpm_qm}\\
  E_{\text{crack}} &\approx 3.35\times10^{-19}\,\text{J} \label{eq:ecrack2}\\
  M_{\text{atomic}} &= M_0(1 - e^{-n_{\text{grad}}/10})\cdot Z \label{eq:mass2}
\end{align}

\noindent\textbf{Complete chain in one line:}
\begin{equation*}
0[\text{vac}] \to \nabla(\text{UA}) \to \text{DPM} \to \mu_s \to U_{g1} \to
U_g\text{-family} \to F_U \to \text{crossing} \to M \to GM/r^2
\end{equation*}
Crossing precedes mass.  Mass is born AT the crossing.  $GM/r^2$ is step 10.

%=============================================================================
\section{Mass Is Dead Until Motion}
\label{sec:dead_mass}
%=============================================================================

When $M_{\text{atomic}}$ first emerges: $v = 0$, $\omega = 0$, $\nabla P = 0$
$\Rightarrow E_{\text{react}} = 0 \Rightarrow F_U = 0$.

\noindent\textbf{Six multiplications from motion:}
\begin{enumerate}[noitemsep]
  \item $E_{\text{react}}$ switches on: $v>0 \Rightarrow 10^{15}\times$ all $U_g$ terms.
  \item $U_{g2}$ shell expands: $(1+\delta_{\text{sw}}v_{\text{sw}})$ factor.
  \item $U_{g3}$ disk activates: $\cos(\omega_s t \pi)$ oscillates.
  \item Buoyancy becomes directional (origin of inertia).
  \item 26-layer cascade: $6{,}279\times$ multiplier from dead-0 to full activation.
  \item Quantum cycle $t_n$ breathes: $U_{g3}$ breathing $=$ EM radiation;
    $U_{g1,g4}$ breathing $=$ gravitational waves.
\end{enumerate}

%=============================================================================
\section{Pre-Mass Roles of Buoyancy, Resonance, Superconductive}
\label{sec:premass}
%=============================================================================

\noindent\textbf{Buoyancy} (the gate):
$U_{b,\text{vac}}(t) = \beta_i \rho_{\text{vac,UA}} \Omega_{\text{vac}}
  \cdot \nabla(\text{UA})/r_{\text{node}} \cdot \cos(\pi t_n)$.
Gate: $E_{\text{crack}} > U_{b,\text{vac}} \Rightarrow$ condenses;
else $\Psi_{\text{proto}}$ disperses.

\noindent\textbf{Resonance} (frequency selector):
$\omega_{\text{DPM}} = n\omega_{\text{vac},n}$ (integer multiples only).
Non-resonant vortices destructively interfere back to vacuum.
This is why atoms have discrete mass values.

\noindent\textbf{Superconductive} (makes mass formation thermodynamically possible):
$g_{\text{SC,pre}} = \sum_{j=1}^4 k_j \rho_{\text{vac,SCm}} H_{\text{SCm}}^{n_j} \Omega_{\text{vac}}^j$.
$H_{\text{SCm}} = 0.99$ persists at stellar core temperatures of $10^6$~K.

%=============================================================================
\section{Proto-Hydrogen and Proto-Helium Evolution}
\label{sec:nucleosynthesis}
%=============================================================================

\noindent\textbf{Proto-Hydrogen ($Z=1$, DPM monomer):}
Single DPM vortex at fundamental resonance.
$M_{\text{proton}} = M_0(1-e^{-1/10}) \approx 0.095\,M_0$.
H-1 dominates ($\sim$75\% by mass) because single-DPM events dominate statistically.

\noindent\textbf{Proto-Helium ($Z=2$, DPM dimer):}
Two vortices in resonance lock: $\omega_{\text{DPM},1} = 2\omega_{\text{DPM},2}$.
He-4 $=$ 4 DPM vortex units (2 protons + 2 neutrons at 90$^\circ$ $U_{g3}$ rotation).
$\sim$25\% by mass (3:1 H:He ratio $=$ DPM density statistic at primordial SCm density).

\noindent\textbf{Nucleosynthesis beyond He:}
$Z \geq 3$ requires $3+$ DPM vortices in simultaneous resonance lock,
requiring stellar $U_{g3}$ disk conditions
($\rho_{\text{SCm}} \sim 10^{15}$~kg/m$^3$).

\noindent\textbf{Dark Hydrogen/Helium:}
SCm itself forms a ``dark'' analogue: $Z_{\text{dark}}=1$ (single SCm vortex, $Q_s=0$).
Their gravitational contribution IS what SM calls dark matter---not separate matter,
but the $U_g$ family above Layer 1.

%=============================================================================
\section{The Sun --- Canonical Solved Example}
\label{sec:sun}
%=============================================================================

\begin{table}[H]
\centering
\begin{tabular}{@{}lll@{}}
\toprule
\textbf{Parameter} & \textbf{Value} & \textbf{Description} \\
\midrule
$M_s$ & $1.989\times10^{30}$ kg & Solar mass \\
$R_s$ & $6.96\times10^8$ m & Solar radius \\
$B_s(t)$ & $(10^{-4} + 0.4\sin\omega_c t)$ T & Surface field, 11-yr cycle \\
$\omega_c$ & $2\pi/3.96\times10^8$ s$^{-1}$ & Solar cycle frequency \\
$\omega_s$ & $2.5\times10^{-6}$ rad/s & Average differential rotation \\
$v_{\text{sw}}$ & $5\times10^5$ m/s & Solar wind speed \\
$R_b$ & $1.496\times10^{13}$ m & Heliosphere radius ($\approx$100 AU) \\
$M_{\text{bh}}$ & $8.15\times10^{36}$ kg & Sgr A* mass \\
$d_g$ & $2.55\times10^{20}$ m & Distance to galactic centre \\
$\Omega_g$ & $7.3\times10^{-16}$ rad/s & Galactic rotation \\
\bottomrule
\end{tabular}
\caption{Solar input parameters for canonical UQFF example.}
\label{tab:solar_params}
\end{table}

\noindent\textbf{Assembled $F_U$ for the Sun:}
\begin{align}
  F_{U,\odot} &=
    (3.38\times10^{16} + 1.35\times10^{19}\sin\omega_c t)
    e^{-0.001t}\cos(\pi t_n) \notag\\
  &+ 8.91\times10^6 \cdot E_{\text{react}} \notag\\
  &+ (10^{-3} + 0.4\sin\omega_c t)\cos(\omega_s t\pi) E_{\text{react}} \notag\\
  &+ (2.26\times10^7 + 9.04\times10^9\sin\omega_c t)
    (1-e^{-5\times10^{-5}t\cos(\pi t_n)}) \times 10^9 E_{\text{react}}
  \label{eq:fu_sun}
\end{align}

\noindent\textbf{$U_{g3}$ counter-rotation generator:}
Equatorial CCW: $\omega_{s,\text{eq}} = 2.9\times10^{-6}$ rad/s.
Coronal CW: $\omega_{s,\text{pol}} = 2.1\times10^{-6}$ rad/s.
The shear layer between these IS the $U_{g3}$ generator.
Without counter-rotation there is no $U_{g3}$ disk.

%=============================================================================
\section{Calibrated Constants}
\label{sec:constants}
%=============================================================================

\begin{table}[H]
\centering
\begin{tabular}{@{}lll@{}}
\toprule
\textbf{Constant} & \textbf{Value} & \textbf{Description} \\
\midrule
$\kappa$ & $5\times10^{-4}$ day$^{-1}$ & $E_{\text{react}}$ exponential decay rate \\
$[\text{SSq}]$ & $0.57$ & Self-similar quotient \\
$\beta_i$ & $0.603$ & Buoyancy coupling constant \\
$H_{\text{SCm}}$ & $0.99$ & SCm superconductive coherence factor \\
$\rho_{\text{vac,SCm}}$ & $7.09\times10^{-37}$ J/m$^3$ & SCm vacuum energy density \\
$\rho_{\text{vac,UA}}$ & $7.09\times10^{-36}$ J/m$^3$ & UA vacuum energy density \\
$v_{\text{SCm}}$ & $c/3 = 10^8$ m/s & SCm velocity (max attraction) \\
$E_{\text{react,0}}$ & $10^{46}$ W/m$^3$ & Astrophysical reactor efficiency base \\
$\alpha$ & $10^{-3}$ day$^{-1}$ & $U_{g1}$ temporal decay rate \\
$E_0$ & $10^{-20}$ J & Base quantum energy (26-level scale) \\
$\eta$ & $10^{-22}$ & Aether coupling constant \\
$\gamma$ & $5\times10^{-5}$ day$^{-1}$ & $U_m$ reciprocation decay rate \\
\bottomrule
\end{tabular}
\caption{Canonical calibrated UQFF constants.  Do not change.}
\label{tab:constants}
\end{table}

%=============================================================================
\section{Millennium Prize Connections}
\label{sec:millennium}
%=============================================================================

\noindent\textbf{Navier--Stokes:}
Quasar jet streams are SCm fluid dynamics against unbound UA.
The TRZ term $\cos(\pi t_n)$ prevents finite-time blowup by enforcing periodic
reversal.  The UQFF fluid element:
$F_{\text{SCm}} = (\rho_{\text{SCm}} v_{\text{SCm}}^2 / r) e^{-\kappa t}$.

\noindent\textbf{Yang--Mills Mass Gap:}
$E_{\text{gap}} = E_{\text{crack}} \approx 3.35\times10^{-19}$~J.
Below this threshold no DPM vortex is sustainable.  The gap is the SCm/UA
attraction energy required to hold a DPM vortex open against vacuum dispersion.

\noindent\textbf{Riemann Hypothesis:}
The 26-layer triadic has natural Riemann structure:
$Z = \text{Li}_{26}([\text{SSq}]) \approx 0.507$.

\noindent\textbf{P $\neq$ NP:}
$\operatorname*{argmin}_n |G^n_{\text{inside}} - O^n_{\text{outside}}|$
requires solving two exponential-time graph problems simultaneously.
Cannot be compressed to polynomial time.

%=============================================================================
\section{Star Magic in Action}
\label{sec:star_magic}
%=============================================================================

\subsection{Heliosphere as Active Synthesis Reactor}

The heliosphere is not a passive boundary.  $U_{g2}$ actively synthesizes solar
wind into hydrogen complexes magnetically adhered to the outer shell.
Heliosphere thickness is a direct measure of stellar age.
Planetary liquid volumes correlate with distance and solar wind exposure.

\subsection{Quasars --- Stellar Reactor Failure Sequence}

$U_{g1}$ coherence degrades $\to$ $U_{g2}$ synthesis destabilizes
$\to$ $U_{g3}$ disk frequency collapses $\to$ SCm expelled at $v_{\text{SCm}} = c/3$
$\to$ expelled SCm ignites against unbound UA $\to$ quasar jets.

Testable prediction: jet luminosity ratio $\propto \text{TRZ}^2$ at each time step.

\subsection{Planetary SCm Inventory}

Saturn's rings $=$ $U_{g3}$ disk made visible by resonance-trapped ice particles.
Uranus's extreme axial tilt $=$ TRZ condition at formation.
Triton's retrograde orbit $=$ late SCm capture event.

%=============================================================================
\section{Implications}
\label{sec:implications}
%=============================================================================

\noindent\textbf{Energy:}
$\gamma = 5\times10^{-5}$ day$^{-1}$; half-life $\approx 55$ years.
$E_{\text{crack}} \approx 3.35\times10^{-19}$~J is within reach of current
high-field laboratory capabilities.

\noindent\textbf{Space Travel:}
A spacecraft with a significant SCm-dense core remains within $U_{g3}$ coupling
range for the entire inner solar system.  Beyond the heliosphere, a vessel with
its own DPM source can maintain orbital coherence without solar $U_{g3}$ coupling.

\noindent\textbf{Quantum Gravity:}
$E_{\text{gap}} = E_{\text{crack}}$.  The gap is not abstract.  It is physical.

\noindent\textbf{Quasar Precursors (Testable):}
Increasing $\delta_{\text{def}}$ amplitude; periodicity shift; heliosphere boundary
fluctuations; planetary orbital resonance drift.  All measurable with current
telescopes.

%=============================================================================
\section*{Conclusion}
\addcontentsline{toc}{section}{Conclusion}
\label{sec:conclusion}
%=============================================================================

The legacy of Star Magic is a unified theory for both quantum and cosmic scales,
built from one foundational mechanism: the maximum attraction of UA and SCm
creating the DPM vortex, which seeds all gravity, which enables all mass.

From the vacuum before the BigBang to quasar jets; from the Yang--Mills mass gap to
the sunspot cycle---all traces back to:
\begin{equation*}
  \text{UA} + \text{SCm} \;\to\; \text{MAX ATTRACTION} \;\to\; \text{DPM}
  \;\to\; F_{U_{Bi}} \text{ REPELS } F_{U_{Bi,i}}
  \;\to\; \text{MATTER}
\end{equation*}

The 26D omega family $\{\omega_i\}_{i=1}^{26}$ is the universal architecture of every
DPM in existence, from the primordial belly button to the smallest atomic nucleus.
Atomic shells, quark flavors, electron charge, electron spin, nuclear magic numbers:
all emerge from the single condition
$\omega_i(r_{i,i+1}) = \omega_{i+1}(r_{i,i+1})$.

Nothing in UQFF is sequential.  Everything is simultaneous.
$GM/r^2$ is the last thing that appears.  Not the first.

%=============================================================================
\appendix
\section{Historical Context}
\label{sec:history}
%=============================================================================

\noindent\textbf{Newton (1687):} $F = GM/r^2$.  Useful observational approximation.
Not the mechanism.  Not the foundation.

\noindent\textbf{Einstein (1915):} GR replaced force with spacetime curvature.
Still began with mass as primary.  Spent his last 30 years searching for the unified
field.  The DPM is what he was searching for.

\noindent\textbf{Standard Model (1920s--present):} Describes particle behaviour
precisely.  Structurally cannot include gravity.  Dark matter, dark energy, quark
confinement, hierarchy problem = symptoms of the missing DPM mechanism.

\noindent\textbf{BCS Superconductivity (1957):} Cooper pairs below critical
temperature.  SCm superconductivity (UQFF): intrinsically superconductive at all
temperatures because $\rho_{\text{SCm}} \sim 10^{15}$~kg/m$^3$ suppresses thermal
decoherence.  $H_{\text{SCm}} = 0.99$ persists at $10^6$~K stellar core.

\bigskip
\begin{center}
\rule{0.5\textwidth}{0.4pt}\\[6pt]
{\small \textcopyright~2025--2026 Daniel T.\ Murphy.
All IP materials are watermarked.  Unauthorized reproduction prohibited.}\\
{\small \texttt{daniel.murphy00@gmail.com}}
\end{center}

\end{document}
